\begin{frame}{왜 마진은 $\frac{2}{\|w\|}$인가 (2/2)}
  \begin{itemize}
    \item 양의 서포트 벡터 $x_+$: $w^\top x_+ + b = 1$ (제약이 등호)
    \item 경계 위의 점 $x_0$: $w^\top x_0 + b = 0$
    \item 두 값을 빼면 $w^\top(x_+ - x_0) = 1$
    \item 위에서 본 대로 이 값이 1만큼 변하려면 $w$ 방향으로
      $\frac{1}{\|w\|}$ 이동 $\Rightarrow$ \textbf{경계에서 서포트
      벡터까지의 직선 거리} $= \frac{1}{\|w\|}$
  \end{itemize}
  \vskip0.6em
  \[
  \text{마진} \;=\; \frac{1}{\|w\|} + \frac{1}{\|w\|}
  \;=\; \frac{2}{\|w\|}
  \]
  \vskip0.4em
  \begin{block}{기하학적 그림이 완성된다}
    $\|w\|$가 작을수록 ``마진 1''을 얻으려면 더 먼 거리까지 떨어져야
    하므로 데이터가 넉넉한 여유를 얻는다.
    \textbf{$\|w\|$를 작게 하는 것 = 마진을 넓게 하는 것}.
    로지스틱회귀의 $w$가 ``확률을 얼마나 급하게 바꾸는가''였다면,
    SVM의 $w$는 ``마진이 얼마나 좁은가''의 역수.
  \end{block}
\end{frame}
